\documentclass{standalone}
\usepackage{standalone}

\begin{document}
\chapter{Implementation}

\section{Technology Stack}
\begin{table}[htbp]
\centering
\caption{Implemented technology stack}
\label{tab:stack}
\small
\begin{tabularx}{\textwidth}{@{}p{2.6cm}p{4.4cm}X@{}}
\toprule
\textbf{Layer} & \textbf{Technology} & \textbf{Role in Sarabo} \\
\midrule
Frontend & React 19, Vite 7 & Component-based single-page app and build tooling \\
Routing & React Router 7 & Public, authenticated and role-specific routes \\
Server state & TanStack Query 5, TanStack Table & Caching, mutations, invalidation; admin data tables \\
UI & Tailwind CSS 4, Radix UI, Motion, Recharts, Leaflet & Design system, accessible primitives, animation, charts, service-area map \\
Forms & react-hook-form & Validated multi-step forms \\
Backend & Node.js, Express 5 & REST API, middleware, controllers, services \\
Database & MongoDB 7 driver, MongoDB Atlas & Persistence and multi-document transactions \\
Identity & Firebase Authentication + Admin SDK & Sign-in (email, Google) and token verification \\
Payments & Stripe 20 & Checkout, signed webhooks, refunds \\
Storage & Supabase Storage & Private images with signed URLs \\
AI & Google Gemini (\texttt{@google/genai}) & Support chat, request helper, offer note, report summary \\
Email & Nodemailer (Gmail) & Request, reminder and job-offer emails \\
Documents & pdfkit & PDF payment receipts \\
Testing & Custom Node test harness, Playwright & Server integration tests; browser tests \\
Delivery & GitHub Actions, Vercel & CI, serverless API (Mumbai region), cron, static hosting \\
\bottomrule
\end{tabularx}
\end{table}

\section{Development Process}
I built Sarabo in small phases that were each easy to review, from 26 July to 4 October 2026, followed by a final improvement phase from 11 to 16 October. Altogether this came to 163 commits on the client and 120 on the server. Every phase followed the same loop. First I defined the rule, then implemented it on the server together with its tests, then built the interface, ran lint, the build and the tests, and finally reviewed the result. Larger features were developed on their own branches and merged through a pull request once CI had passed (Table~\ref{tab:phases}).

\begin{longtable}{@{}p{3.2cm}p{2.2cm}p{8.8cm}@{}}
\caption{Development phases (from the git history)}\label{tab:phases}\\
\toprule
\textbf{Phase} & \textbf{Dates} & \textbf{Main work} \\
\midrule
\endfirsthead
\toprule
\textbf{Phase} & \textbf{Dates} & \textbf{Main work} \\
\midrule
\endhead
\bottomrule
\endfoot
1. Foundation and hardening & 26--28 Jul & Imported the base codebase; secured authentication and role routing; secure Stripe checkout and webhook; validated status transitions; public sanitized tracking; transactional assignment, approval and completion. \\
2. Admin and notifications & 29 Jul -- 4 Aug & Admin request management; in-app notifications; production database isolation; health monitoring; double-booking prevention. \\
3. Repair domain (V2) & 5--8 Aug & Service taxonomy and priced catalogue; technician expertise; V2 request schema; expertise-based matching; damage uploads; prices in BDT; inspection, quote, payment and repair workflows; safe deletion; verified email. \\
4. Rename to Sarabo & 10--12 Aug & Migrated database, collections, fields, modules, API routes and UI from parcel to repair vocabulary; assignment accept/reject; receipt confirmation. \\
5. Money and UI redesign & 13--19 Aug & 90/10 settlement; wallet and withdrawals; storage moved to Supabase; complete ``service spine'' redesign; accessibility pass. \textit{(Earlier report submitted 20 Aug.)} \\
6. AI support and feedback & 18--24 Sep & Gemini support chat; reviews and private reports with moderation; product tour; second redesign; admin catalogue editor; CI on GitHub Actions. \\
7. Operations & 24--25 Sep & Pickup scheduling; region rule; email updates; Playwright tests; admin reports and attention list; in-app reminders; refusing unmatchable requests. \\
8. Job portal and inspection fee & 25 Sep -- 2 Oct & Technician offers and customer choice; inspection-fee payment, credit, refunds and split; reject-and-repost; public jobs board; job notifications; PDF receipts; one-step approve-and-pay. \\
9. AI helpers and payouts & 2--4 Oct & AI request helper (text and photo); price guide and AI offer note; AI report summary; bKash withdrawals; hosting configuration. \\
10. Final improvements & 11--16 Oct & Map search, final payout on removal, faster API and AI, offer ranking, warranty, deadlines for every waiting state, pickup condition, minimum inspection fee, technician profile editing (Section~\ref{sec:finalphase}). \\
\end{longtable}

\section{Frontend Implementation}
\subsection{Role-Aware Routing}
The dashboard is behind a route that requires sign-in, with guards for each role. Customers see Overview, Create Request, My Requests and Payment History. Technicians see Overview, Find jobs, Assigned jobs, Completed jobs, Warranty claims and Wallet. Admins see Overview, Repair requests, Assign technicians, Technicians, Withdrawals, Warranty claims, Users, Services \& prices, Reports, Technician reports and Technician reviews. Every role also has Notifications and Profile, and on phones these pages move to a tab bar at the bottom. Deep links such as \texttt{?category=smartphone} are kept through sign-in. The routing is only there for convenience, though. Every protected endpoint checks authentication and authorization again on the server.

\subsection{Server-State Hooks}
Each feature has its own query-key factory (for example \texttt{quoteKeys} or \texttt{walletKeys}) and hooks for reading and changing data. When something changes, only the queries that are now out of date are refreshed. Approving a quote, for instance, refreshes the quote, the request, the request list and the payment eligibility. This keeps what all three roles see in step, without anyone having to reload the page.

\subsection{Action-First Workspace}
Every request page (see Figure~\ref{fig:wt-c10} for an example) starts with the service spine. Below it is a ``next step'' panel that tells the signed-in user what to do, such as ``Choose a technician'' or ``Approve and pay BDT 3,800'', and the supporting details come after that. Shared components for empty, loading, error and confirmation states and for forms make every page behave the same way. A product tour introduces each dashboard, and a command menu lets users move around with the keyboard.

\section{Backend Implementation}
\subsection{Controllers, Services and Adapters}
Each domain has a controller for its endpoints, for example \texttt{jobPortalController}, \texttt{quoteController} and \texttt{paymentController}. Logic that is shared goes into services such as \texttt{jobBooking}, \texttt{inspectionFeeRefunds}, \texttt{paymentEligibility}, \texttt{technicianMatchGate} and \texttt{notificationService}. The external providers (Supabase, Gemini and Stripe) are wrapped in adapters, which the tests replace with fakes.

\subsection{Database Transactions}
Whenever several records depend on each other, they are updated inside a MongoDB transaction. Take booking after an inspection-fee payment. In one transaction (\texttt{services/\allowbreak jobBooking.js}) the server records the payment, sets the request to \texttt{driver\_assigned}, marks the technician as busy, accepts the chosen offer and marks all the other offers \texttt{not\_chosen}. If any of these steps fails, none of them is saved.

\subsection{Notifications, Email and Scheduled Jobs}
All in-app notifications are created through one service. Each one has a de-duplication key, so retrying an action never sends the same notification twice. Emails (request received, pickup reminder and new job) are sent through Gmail after the main action has finished. If an email fails, it is logged in \texttt{email\_log}, and the request itself still succeeds.

Once a day a Vercel Cron job calls \texttt{GET /cron/pickup-reminders} at 12:00 UTC, which is 18:00 in Bangladesh. The call is protected by a secret. The job sends the pickup reminders for the next day and applies every deadline described in Section~\ref{sec:nostuck}. It closes declined and unpaid quotes that nobody acted on, confirms receipts that the customer has left unanswered for three days, and warns admins about late warranty steps and stalled work. Each of these steps runs separately, so if one fails the others still run. Email and the daily job are both turned on in production, which means the live site sends reminders by email as well as in the app.

\section{Service Catalogue and Matching}
Admins edit services and prices from the dashboard. Changing a price increases \texttt{pricingRule.version}, while requests that already exist keep their own copy of the old price. Every service has an inspection fee of at least 100\,BDT, or 200\,BDT if the top of its estimate is above 1,000\,BDT. A single rule, \texttt{minimumInspectionFee}, is used for requests, offers, the catalogue editor and the public price list. A one-time migration raised the fee for the 16 live services that had none.

When a request is created, \texttt{technicianMatchGate} checks that there is an approved technician in the region who covers the device and the repair at the required level. If there is not, the API answers \texttt{409 NO\_TECHNICIAN\_AVAILABLE}, and the refused request is counted once per customer, service, region and day, with the customer stored only as a hash. These counts show admins where they need to recruit more technicians.

\section{Pickup Scheduling}
Customers choose one of four two-hour windows between 10:00 and 18:00 Bangladesh time, from Saturday to Thursday, anywhere from one hour to seven days ahead. Each region accepts a set number of pickups per slot. This number can be configured and is three by default, and a full slot returns \texttt{409 PICKUP\_SLOT\_FULL}. The customer can change the window until the device is collected, and if the pickup moves, the reminder is sent again.

\section{Quotation Model}
A quote is made up of labour, parts and additional charges, all in whole taka, plus a total, a currency, a version and a status. When a quote is rejected, it is moved to \texttt{quoteHistory} before the live quote is cleared. That way there is still a record of it, and it can never be paid by mistake.

\begin{lstlisting}[caption={Quote total and amount still to pay},label={lst:quote}]
repairSubtotal   = laborAmount + partsAmount + additionalCharges
quote.totalAmount = repairSubtotal            (currency "BDT")
amountDue        = max(quote.totalAmount - inspectionFeePaid, 0)
\end{lstlisting}

\section{Payment Implementation}
Every Stripe Checkout session is created from data the server has stored, either the approved quote or the fee of the chosen offer. It never uses an amount sent by the browser. The session's metadata records what it is paying for, either \texttt{repair} or \texttt{inspection\_fee}. The signed webhook and, as a backup, the success page both call the same idempotent processor, so each payment is recorded exactly once. A single server function decides whether a request counts as paid, and the eligibility checks, admin filters and status badges all use it. Every payment also has a PDF receipt showing the Stripe ID, the date in Dhaka time, the device, the repair, the technician, the amount and the refund status.

\section{Technician Settlement and Wallet}
As soon as a repair is paid, the server takes a snapshot of the settlement from the approved quote, and this snapshot does not change afterwards:

\begin{lstlisting}[caption={Settlement rule (whole taka)},label={lst:settlement}]
commission  = round(repairSubtotal * 10 / 100)    # Sarabo
receivable  = repairSubtotal - commission          # Technician
# invariant: commission + receivable == repairSubtotal
\end{lstlisting}

The settlement stays \emph{pending} until the customer confirms they have received the device, and then it becomes \emph{available}. Balances are never stored. They are worked out each time from the settlements and withdrawals, so they cannot slowly drift away from the truth. A technician can withdraw up to their available balance to a bKash number. An admin then sends the money with bKash Send Money and records the TrxID. A unique partial index makes sure each technician can only have one open withdrawal at a time.

\section{Damage and Evidence Images}
Images go straight from the browser to a private Supabase bucket through a signed upload URL. They must be JPEG, PNG or WebP files of up to 8\,MB, with at most three per request. Before attaching an image, the server checks the real type and size of the stored file. Signed links for viewing it last five minutes and are given only to the owner, an admin or the assigned technician. When something is deleted, the storage clean-up is retried until it succeeds.

\section{Public Tracking and Jobs Board}
Every request gets a tracking code of the form \texttt{SRB-<random>}, with 128 bits of randomness. \texttt{GET /public/trackings/:code} only accepts exact codes, is rate limited, and returns nothing but the status timeline. The jobs board (\texttt{GET /jobs}, limited to 60 requests per minute per address) shows the device, the repair, the area, the pickup day, the usual price and the number of offers. It never shows the problem description, photos, the tracking code or any personal details.

\section{Deployment and Continuous Integration}
The API is a Vercel serverless function (Node.js 24) at \url{https://sarabo-server.vercel.app}, and the client is a static Vercel site at \url{https://sarabo-client-six.vercel.app}. Both are deployed from the \texttt{main} branch. The API runs in Vercel's Mumbai region (\texttt{bom1}), close to the MongoDB Atlas cluster. Moving it there brought the first request after an idle period down from 10.7\,s to 2.9\,s, and a normal request from 0.75\,s to about 0.2\,s. Production has a database of its own, and the server refuses to start there unless its secrets and an \texttt{https} client origin are set.

GitHub Actions runs on every push. It runs the server suites on a throwaway MongoDB replica set, never with real credentials, and it runs lint, the build and the browser tests for the client. I only push to \texttt{main} after all of these checks have passed.

\section{Final Improvement Phase}
\label{sec:finalphase}
The features were finished on 4 October. After that, a final phase from 11 to 16 October 2026 got the project ready for submission. It began with a review of both codebases and a fresh run of every check (Chapter~\ref{ch:testing}). I recorded each change in the project work log, with the problem it solved, the files it touched and how it was verified. Table~\ref{tab:finalphase} summarises these changes.

{\small
\begin{longtable}{@{}p{0.6cm}p{2cm}p{11.4cm}@{}}
\caption{Final improvement phase log}\label{tab:finalphase}\\
\toprule
\textbf{\#} & \textbf{Date} & \textbf{Change and verification} \\
\midrule
\endfirsthead
\toprule
\textbf{\#} & \textbf{Date} & \textbf{Change and verification} \\
\midrule
\endhead
\bottomrule
\endfoot
0 & 11 Oct 2026 & Codebase and documentation review; re-ran server suites against an isolated in-memory MongoDB replica set, client lint, production build and Playwright browser tests. \\
1 & 11 Oct 2026 & Service-area map: OpenStreetMap coordinates for 226 of the 236 covered sub-areas, each kept only if it lies inside its own district. Searching a sub-area (e.g.\ Ambarkhana) now zooms to that spot instead of the district centre. Verified with two new browser tests on desktop and mobile. \\
2 & 11 Oct 2026 & Final payout on technician removal: removing a technician queues their whole available balance as a bKash withdrawal in the same transaction (Section~\ref{sec:removalpayout}). Verified with 15 new server checks. \\
3 & 11 Oct 2026 & Homepage hero: the photo moved to the right side (top band on phones) so the headline no longer sits over it. Checked with screenshots at three screen widths. \\
4 & 11 Oct 2026 & Final verification of the committed code and CI on both repositories (Section~\ref{sec:results}). \\
5 & 14 Oct 2026 & Slow first load: the API was moved to Vercel's Mumbai region, next to the database (10.7\,s $\rightarrow$ 2.9\,s cold, 0.75\,s $\rightarrow$ 0.2\,s warm, measured live). \\
6 & 14 Oct 2026 & AI support timeouts: faster Gemini model with an automatic fallback model; live answers now take 1.4--1.9\,s. 4 new tests. \\
7 & 14 Oct 2026 & Smarter matching: offers ranked by rating, reliability, price, district and reply time, with a ``Best match'' badge and reasons (Section~\ref{sec:ranking}). 22 new tests. \\
8 & 14 Oct 2026 & 30-day repair warranty with free re-inspection, re-repair, admin review and refunds (Section~\ref{sec:warranty}). 45 new tests. \\
9 & 14 Oct 2026 & Whole-workflow review, then fixes so no request can get stuck: unpaid quotes, not-repairable devices, repairs that cannot be finished, unconfirmed receipt, warranty deadlines (Section~\ref{sec:nostuck}). 42 new tests. \\
10 & 14 Oct 2026 & Remaining gaps: late payment after cancellation refunded, device return after cancellation, stalled-work alerts, pickup condition, customer not home. 30 new tests. \\
11 & 14 Oct 2026 & Photos upload only when the customer presses ``Upload''; technicians edit their own profile; minimum inspection fee for every repair, applied to old offers and migrated on production. 22 new tests. \\
12 & 14 Oct 2026 & All of the above pushed and deployed; full verification (Section~\ref{sec:results}) and screenshots of the running system (Chapter~\ref{ch:walkthrough}). \\
13 & 16 Oct 2026 & Stripe's repair-payment page named the item ``Repair request: undefined''; it now names the device (``Repair: Samsung Galaxy A54''). 1 new test; pushed and deployed. \\
\end{longtable}
}

\end{document}
